\section{Analisis Solusi}
\label{sec:analisis-solusi}

Berdasarkan masalah yang telah diidentifikasi pada Subbab \ref{subsec:identifikasi-masalah}, solusi yang diperlukan bukan hanya penggantian format token, tetapi pembaruan cara DecMed mengekspresikan, menegakkan, dan mendelegasikan hak akses. DecMed sudah memiliki fondasi yang sesuai untuk sistem RME terdesentralisasi melalui IOTA, IPFS, dan PRE. Oleh karena itu, solusi yang dianalisis pada bagian ini diarahkan untuk memperkuat lapisan kontrol akses tanpa mengubah peran utama komponen-komponen tersebut.

Tiga kebutuhan utama yang menjadi dasar solusi adalah kontrol akses granular, segmentasi data \textit{off-chain}, dan delegasi berantai dengan atenuasi \textit{offline}. Ketiganya saling bergantung. Token yang granular tidak cukup apabila data yang dilindungi masih berbentuk objek terenkripsi yang terlalu besar. Sebaliknya, segmentasi data tidak dapat dimanfaatkan secara aman apabila token tidak mampu menyatakan batasan akses secara eksplisit. Delegasi berantai juga memerlukan token yang dapat dipersempit tanpa membuka peluang perluasan hak akses.

\subsection{Analisis Perbandingan Token Otorisasi}
\label{subsec:analisis-perbandingan-token-otorisasi}

Berdasarkan masalah yang telah diidentifikasi pada Subbab \ref{subsec:identifikasi-masalah}, mekanisme otorisasi pada DecMed perlu mendukung kontrol akses granular, pembatasan cakupan akses, delegasi hak akses yang dapat dipersempit, verifikasi token secara mandiri, serta kesesuaian dengan arsitektur terdistribusi yang melibatkan IOTA, IPFS, dan PRE. Oleh karena itu, diperlukan analisis perbandingan terhadap beberapa pendekatan token otorisasi untuk menentukan mekanisme yang paling sesuai dengan kebutuhan sistem.

Alternatif token otorisasi yang dibandingkan pada analisis ini adalah JWT, PASETO, Macaroons, dan Biscuit. JWT dipilih sebagai pembanding karena merupakan bentuk token yang digunakan pada rancangan DecMed awal dan umum digunakan untuk merepresentasikan klaim otorisasi. PASETO dipilih karena memiliki tujuan penggunaan yang serupa dengan JWT sebagai token \textit{stateless}, tetapi dengan desain kriptografi yang lebih terarah dan lebih membatasi pilihan algoritma untuk mengurangi risiko kesalahan konfigurasi. Macaroons dianalisis sebagai kandidat utama karena menyediakan mekanisme \textit{caveat} dan atenuasi hak akses secara lokal, sehingga mendukung delegasi terbatas tanpa mengharuskan penerbit token membuat token baru. Biscuit dianalisis sebagai pembanding karena juga mendukung verifikasi terdesentralisasi dan atenuasi \textit{offline}, tetapi menggunakan pendekatan yang lebih ekspresif melalui aturan berbasis \textit{logic language}.

Kriteria evaluasi pada Tabel \ref{tab:perbandingan-token-otorisasi} diturunkan dari kebutuhan otorisasi DecMed. Kemampuan membawa klaim atau batasan akses granular digunakan untuk menilai sejauh mana token dapat merepresentasikan informasi otorisasi, seperti identitas subjek, pasien terkait, jenis data, fungsi klinis, aksi yang diperbolehkan, dan masa berlaku akses. Delegasi hak akses yang dipersempit dan kemampuan atenuasi \textit{offline} digunakan untuk menilai apakah pemegang token dapat menurunkan sebagian hak akses kepada aktor lain tanpa memperluas cakupan akses awal dan tanpa selalu meminta penerbit token membuat token baru. Ekspresivitas pembatasan akses digunakan untuk menilai kemampuan token dalam menyatakan aturan atau kondisi akses yang lebih spesifik. Verifikasi mandiri oleh komponen penerima diperlukan agar token dapat divalidasi oleh komponen seperti \textit{Server} PRE tanpa selalu bergantung pada server otorisasi pusat. Selain itu, kesesuaian dengan kebutuhan delegasi DecMed dan kompleksitas integrasi juga dipertimbangkan karena mekanisme token yang dipilih harus dapat diterapkan dalam arsitektur DecMed yang melibatkan IOTA, IPFS, PRE, dan aplikasi \textit{client}, tanpa menambah kompleksitas implementasi di luar ruang lingkup penelitian.


\begin{table}[!ht]
	\centering
	\caption{Perbandingan Pendekatan Token Otorisasi}
	\label{tab:perbandingan-token-otorisasi}
	\scriptsize
	\begin{tabular}{|p{0.24\textwidth}|p{0.14\textwidth}|p{0.14\textwidth}|p{0.14\textwidth}|p{0.14\textwidth}|}
		\hline
		\textbf{Kriteria} & \textbf{JWT}                                     & \textbf{PASETO} & \textbf{Macaroons} & \textbf{Biscuit} \\ \hline

		Kemampuan membawa batasan akses
		                  & Melalui klaim
		                  & Melalui klaim
		                  & Melalui \textit{caveats}
		                  & Melalui \textit{fact} dan \textit{rule}                                                                    \\ \hline

		Delegasi hak akses
		                  & Membutuhkan mekanisme tambahan
		                  & Membutuhkan mekanisme tambahan
		                  & Didukung melalui penambahan \textit{caveats}
		                  & Didukung melalui penambahan \textit{rule}                                                                  \\ \hline

		Kemampuan atenuasi \textit{offline}
		                  & Tidak
		                  & Tidak
		                  & Ya
		                  & Ya                                                                                                         \\ \hline

		Ekspresivitas pembatasan akses
		                  & Sedang
		                  & Sedang
		                  & Sedang--tinggi
		                  & Tinggi                                                                                                     \\ \hline

		Verifikasi mandiri oleh komponen penerima
		                  & Ya, dengan kunci verifikasi yang sesuai
		                  & Ya, dengan kunci verifikasi yang sesuai
		                  & Ya, dengan kunci dan pemeriksaan \textit{caveat}
		                  & Ya, dengan kunci publik dan pemeriksaan aturan                                                             \\ \hline

		Kesesuaian dengan kebutuhan delegasi pada DecMed
		                  & Rendah--sedang
		                  & Rendah--sedang
		                  & Tinggi
		                  & Tinggi                                                                                                     \\ \hline

		Kompleksitas integrasi
		                  & Rendah
		                  & Rendah--sedang
		                  & Sedang
		                  & Tinggi                                                                                                     \\ \hline
	\end{tabular}
\end{table}


JWT dan PASETO memiliki kemiripan dari sisi penggunaan sebagai token \textit{stateless} yang membawa klaim. Keduanya dapat digunakan untuk menyatakan informasi otorisasi seperti identitas subjek, cakupan akses, dan waktu kedaluwarsa. Namun, pembatasan tersebut umumnya ditentukan pada saat token diterbitkan. Dengan demikian, apabila penerima token ingin mendelegasikan sebagian hak akses kepada pihak lain, sistem tetap memerlukan penerbitan token baru atau mekanisme kontrol tambahan di sisi \textit{server}. Karakteristik ini kurang sesuai dengan kebutuhan DecMed yang memerlukan delegasi antar aktor layanan kesehatan tanpa selalu melibatkan pasien atau \textit{server} otorisasi sebagai penerbit token baru.

Biscuit menyediakan kemampuan yang lebih dekat dengan kebutuhan DecMed karena mendukung atenuasi \textit{offline} dan aturan otorisasi yang ekspresif. Namun, ekspresivitas Biscuit juga membawa kompleksitas tambahan karena pemeriksaan otorisasi dilakukan menggunakan pendekatan berbasis fakta dan aturan. Dalam konteks penelitian ini, kebutuhan pembatasan akses DecMed masih dapat direpresentasikan melalui batasan yang relatif langsung, seperti identitas pasien, identitas aktor, dataset, fungsi klinis, aksi, masa berlaku, dan batas delegasi. Oleh karena itu, penggunaan Biscuit dinilai lebih kompleks dibandingkan kebutuhan implementasi yang ingin dicapai.

Berdasarkan perbandingan tersebut, Macaroons dipilih sebagai mekanisme token otorisasi pada DecMed karena memberikan keseimbangan antara kemampuan kontrol akses granular, dukungan delegasi terbatas, atenuasi hak akses secara lokal, dan kompleksitas implementasi yang masih sesuai dengan ruang lingkup penelitian. Melalui mekanisme \textit{caveats}, token dapat membawa pembatasan akses seperti identitas pasien, identitas aktor, jenis data, aksi yang diperbolehkan, masa berlaku, serta batasan delegasi. \textit{Caveats} tersebut juga dapat ditambahkan secara lokal oleh pemegang token untuk mempersempit hak akses sebelum token didelegasikan kepada pihak lain, tanpa mengharuskan penerbit token membuat token baru. Karakteristik ini sesuai dengan kebutuhan DecMed untuk mendukung delegasi terbatas dalam alur pelayanan kesehatan, misalnya dari petugas administrasi kepada dokter, atau dari dokter kepada petugas laboratorium, tanpa memperluas cakupan akses yang diberikan oleh pasien. Dibandingkan JWT dan PASETO, Macaroons lebih sesuai karena mendukung penambahan pembatasan akses oleh pemegang token. Dibandingkan Biscuit, Macaroons memiliki ekspresivitas yang lebih sederhana, tetapi sudah mencukupi untuk merepresentasikan kebutuhan pembatasan akses DecMed melalui \textit{caveats}. Dengan demikian, Macaroons dipilih bukan sebagai token yang paling unggul untuk seluruh skenario otorisasi, melainkan sebagai mekanisme yang paling selaras dengan kebutuhan khusus DecMed.
\subsection{Analisis Segmentasi Data}
\label{subsec:analisis-segmentasi-data}

Masalah granularitas tidak hanya berada pada token, tetapi juga pada struktur data yang dilindungi. Jika data RME disimpan sebagai satu \textit{encrypted blob}, maka pemberian akses terhadap blob tersebut akan membuka seluruh isi data meskipun aktor hanya membutuhkan sebagian kecil informasi. Oleh karena itu, implementasi pemrosesan data pada DecMed perlu disesuaikan untuk memenuhi kebutuhan akses granular sebagaimana dijelaskan pada Subbab \ref{subsec:identifikasi-masalah}.

Alternatif yang dianalisis adalah melakukan segmentasi data sejak proses penyimpanan \textit{off-chain}. Pada pendekatan ini, data klinis dipisahkan menjadi beberapa objek berdasarkan kategori tertentu sebelum disimpan pada IPFS. Dengan pendekatan tersebut, setiap kategori data dapat diambil secara independen sehingga akses terhadap data menjadi lebih granular.

Namun, pendekatan tersebut memiliki keterbatasan implementasi. IPFS sebagai media penyimpanan \textit{off-chain} tidak menyediakan mekanisme pencarian atau pemfilteran data secara \textit{native}. Pemisahan data menjadi banyak objek akan membutuhkan sistem indeks tambahan untuk melakukan pencarian dan pengelolaan relasi antar segmen data. Kompleksitas tersebut meningkatkan kebutuhan sinkronisasi metadata dan berada di luar ruang lingkup implementasi penelitian ini.

Berdasarkan analisis tersebut, pendekatan yang dipilih adalah melakukan segmentasi data klinis berdasarkan fungsi klinis dengan memanfaatkan metadata yang disimpan pada \textit{on-chain}. Pada pendekatan ini, data klinis dipisahkan menjadi beberapa segmen sebelum disimpan pada penyimpanan \textit{off-chain}. Setiap segmen kemudian direpresentasikan melalui metadata pada \textit{smart contract} sehingga proses pemilahan data dapat dilakukan sebelum sistem mengambil data dari IPFS.

Pendekatan ini lebih memungkinkan untuk diimplementasikan karena \textit{smart contract} dapat digunakan untuk melakukan pemetaan dan pemfilteran berdasarkan kategori data tanpa memerlukan mekanisme pencarian langsung pada IPFS. Selain itu, proses pengambilan data menjadi lebih terarah karena sistem hanya mengambil segmen data yang sesuai dengan kategori yang diminta.

Penentuan kategori data tidak dilakukan secara arbitrer, melainkan mengacu pada struktur data rekam medis elektronik yang ditetapkan dalam Keputusan Menteri Kesehatan Republik Indonesia Nomor HK.01.07/MENKES/1423/2022 tentang Pedoman Variabel dan Metadata pada Penyelenggaraan Rekam Medis Elektronik \autocite{kemenkes1423-2022}. Pada pedoman tersebut, data rekam medis dibagi menjadi beberapa \textit{dataset} pelayanan, yaitu

\begin{itemize}[]
    \item Instalasi Gawat Darurat (IGD)
    \item Rawat Jalan
    \item Rawat Inap
    \item Laboratorium
    \item Apotek
\end{itemize}
sebagaimana dijelaskan pada Subbab \ref{subsec:isi-rme}.

Pembagian berdasarkan \textit{dataset} pelayanan dipilih karena setiap layanan memiliki alur kerja, jenis data, dan kebutuhan akses yang berbeda. Sebagai contoh, petugas laboratorium hanya memerlukan akses terhadap data pemeriksaan laboratorium, sedangkan tenaga farmasi lebih berfokus pada data terapi dan obat pasien. Dengan demikian, pemisahan berdasarkan \textit{dataset} memberikan batas akses yang lebih sesuai dengan konteks pelayanan medis.

Selain pembagian berdasarkan \textit{dataset}, data klinis juga dianalisis berdasarkan kategori fungsi klinis. Pada \textit{dataset} Rawat Jalan, data administratif seperti Lembar Identitas, Cara Pembayaran, dan \textit{General Consent} tidak diperinci lebih lanjut pada Tugas Akhir ini karena lebih berfokus pada kebutuhan administratif dan tidak berkaitan langsung dengan mekanisme kontrol akses granular terhadap data klinis.

Tugas Akhir ini difokuskan pada kategori fungsi klinis yang memiliki kebutuhan akses berbeda antar tenaga kesehatan. Pada Formulir Umum / Asesmen Awal Rawat Jalan, kategori fungsi klinis yang digunakan meliputi:

\begin{itemize}
    \item Anamnesis
    \item Pemeriksaan Fisik
    \item Pemeriksaan Psikologi
\end{itemize}

Sedangkan pada Pemeriksaan Spesialistik, kategori fungsi klinis yang digunakan meliputi:

\begin{itemize}
    \item Riwayat Penggunaan Obat
    \item Rencana Rawat
    \item Instruksi Medik dan Keperawatan
    \item Pemeriksaan Penunjang
    \item Diagnosis
    \item Persetujuan Tindakan / Penolakan Tindakan (\textit{Informed Consent})
    \item Terapi
\end{itemize}

Melalui pendekatan ini, data \textit{off-chain} tidak lagi diperlakukan sebagai satu objek data besar, melainkan sebagai kumpulan segmen data klinis yang memiliki kategori yang terdefinisi. Segmentasi tersebut memungkinkan penerapan kontrol akses granular karena kebijakan akses dapat dicocokkan dengan metadata kategori yang tersimpan pada \textit{on-chain}. Selain itu, pendekatan ini juga memberikan objek penegakan yang lebih jelas bagi \textit{caveats} pada Macaroons, karena batasan seperti \texttt{dataset\_category} dapat dipetakan secara langsung terhadap kategori data klinis yang diakses.
% \subsection{Analisis Integrasi dengan PRE}
\label{subsec:analisis-integrasi-pre}

PRE tetap menjadi titik penting dalam rancangan karena data RME berada dalam bentuk terenkripsi. Pada DecMed, PRE digunakan untuk mengubah akses dekripsi dari pemilik data ke penerima akses tanpa membuka \textit{plaintext}. Dalam solusi yang diusulkan pada Tugas Akhir ini, PRE tidak hanya menjalankan proses \textit{re-encryption}, tetapi juga memiliki fungsi baru sebagai titik penegakan kebijakan akses.

Sebelum melakukan \textit{re-encryption}, \textit{server} PRE perlu memvalidasi token Macaroon dan mengevaluasi seluruh \textit{caveats}. Jika permintaan akses tidak sesuai dengan batasan token, proses \textit{re-encryption} tidak dilakukan. Dengan cara ini, token Macaroon, metadata IOTA, dan segmen data IPFS bekerja sebagai satu mekanisme kontrol akses yang saling melengkapi.

% ?Kenapa ga bikin instance baru aja? perlu dijelaskan

\subsection{Pemetaan Kebutuhan terhadap Solusi}
\label{subsec:pemetaan-kebutuhan-solusi}

Hasil analisis solusi dirangkum pada Tabel \ref{tab:pemetaan-kebutuhan-solusi}. Pemetaan ini menunjukkan hubungan antara kebutuhan sistem yang ditemukan pada analisis masalah dengan komponen solusi yang akan dirancang pada bagian berikutnya.

\begin{table}[!ht]
	\centering
	\caption{Pemetaan Kebutuhan Sistem terhadap Analisis Solusi}
	\label{tab:pemetaan-kebutuhan-solusi}
	\small
	\begin{tabular}{|p{0.28\textwidth}|p{0.3\textwidth}|p{0.32\textwidth}|}
		\hline
		\textbf{Kebutuhan Sistem}                 & \textbf{Komponen Solusi}                      & \textbf{Implikasi terhadap Rancangan}                                       \\ \hline
		Kontrol akses granular                    & Macaroons dengan \textit{caveats}             & Token perlu memuat batasan aktor, kategori data, hak akses, dan waktu.      \\ \hline
		Segmentasi data \textit{off-chain}        & Pemisahan data IPFS berdasarkan kategori data & Metadata segmen data perlu dapat dicocokkan dengan \textit{caveats} token.  \\ \hline
		Delegasi berantai                         & Atenuasi \textit{offline} pada Macaroons      & Pemegang token dapat menurunkan hak akses dengan cakupan yang lebih sempit. \\ \hline
		Penegakan akses terhadap data terenkripsi & Validasi token pada \textit{server} PRE       & PRE perlu memvalidasi token sebelum menjalankan \textit{re-encryption}.     \\ \hline
	\end{tabular}
\end{table}

Dengan hasil analisis tersebut, solusi yang dirancang pada bagian berikutnya difokuskan pada empat aspek: arsitektur integrasi Macaroons ke DecMed, struktur token dan \textit{caveats}, segmentasi data \textit{off-chain}, serta alur akses dan delegasi.

